\documentclass[11pt]{article}
\usepackage[margin=1in]{geometry}
\usepackage{enumitem}
% =============================================================================
% Preambles/header.tex — shared LaTeX/Beamer preamble for this project
%
% Use from a Beamer lecture by prepending:
%     \input{header}
% and compiling with TEXINPUTS=../Preambles:$TEXINPUTS (the /compile-latex
% skill does this automatically).
%
% PALETTE CONTRACT. Color names defined here MUST match the names in
% Quarto/theme-template.scss so Beamer and Quarto renderings use the same
% palette. The scripts/check-palette-sync.sh script enforces this.
%
% To customize: edit the HEX values below AND the matching $variable in
% Quarto/theme-template.scss, then run scripts/check-palette-sync.sh.
% =============================================================================

% -----------------------------------------------------------------------------
% Engine detection + encoding
%
% Our /compile-latex skill uses XeLaTeX, where inputenc is unnecessary and
% can produce encoding warnings. Only load inputenc under pdfTeX.
% -----------------------------------------------------------------------------
\usepackage{iftex}
\ifPDFTeX
  \usepackage[utf8]{inputenc}
\fi

\usepackage{amsmath, amssymb, amsthm}
\usepackage{booktabs}
\usepackage{graphicx}

% -----------------------------------------------------------------------------
% PALETTE — mirrors Quarto/theme-template.scss variable names.
% Load xcolor and define colors BEFORE anything that references them
% (notably hyperref, hypersetup, and the Beamer color assignments).
% -----------------------------------------------------------------------------
\usepackage{xcolor}

\definecolor{primary-blue}{HTML}{012169}      % headings, accents
\definecolor{primary-gold}{HTML}{B9975B}      % emphasis, borders
\definecolor{highlight-yellow}{HTML}{F2A900}  % markers, alerts
\definecolor{light-bg}{HTML}{E8EDF5}          % title / section backgrounds
\definecolor{jet}{HTML}{1A1A1A}               % body text

% Semantic colors (used in diagrams and annotations)
\definecolor{positive}{HTML}{15803D}          % good / observed / identified
\definecolor{negative}{HTML}{B91C1C}          % bad / problematic / bias
\definecolor{neutral}{HTML}{525252}           % reference / context

% Highlight accents (match .hi-* classes in the SCSS)
\definecolor{hi-slate}{HTML}{314F4F}
\definecolor{hi-green}{HTML}{2E8B57}
\definecolor{hi-red}{HTML}{C41E3A}

% Semantic aliases so skill templates and snippets can write
% \textcolor{accent}{...} without hard-coding "primary-blue".
\colorlet{accent}{primary-blue}
\colorlet{accent2}{primary-gold}

% -----------------------------------------------------------------------------
% Hyperref — loaded AFTER xcolor + palette so link colors resolve.
% -----------------------------------------------------------------------------
\usepackage{hyperref}
\hypersetup{colorlinks=true, linkcolor=primary-blue, urlcolor=primary-blue,
            citecolor=primary-blue, pdfborder={0 0 0}}

% -----------------------------------------------------------------------------
% Course code — set once per deck (after \input{header}, before \begin{document})
% via \coursecode{CS401}. Shown in the footer on every slide so a deck's course
% is identifiable without opening the title page. Empty by default.
% -----------------------------------------------------------------------------
\makeatletter
\newcommand{\coursecode}[1]{\gdef\@coursecodetext{#1}}
\gdef\@coursecodetext{}
\makeatother

% -----------------------------------------------------------------------------
% Beamer theme assignments (only apply under Beamer).
%
% We detect Beamer via \@ifclassloaded{beamer}, which is the documented,
% non-internal way. This must be wrapped in \makeatletter/\makeatother
% because \@ifclassloaded contains an @-catcode character.
% -----------------------------------------------------------------------------
\makeatletter
\@ifclassloaded{beamer}{%
  \usefonttheme{professionalfonts}
  \setbeamercolor{structure}{fg=primary-blue}
  \setbeamercolor{frametitle}{fg=primary-blue}
  \setbeamercolor{title}{fg=primary-blue}
  \setbeamercolor{subtitle}{fg=primary-gold}
  \setbeamercolor{author}{fg=primary-blue}
  \setbeamercolor{institute}{fg=neutral}
  \setbeamercolor{date}{fg=primary-gold}
  \setbeamercolor{itemize item}{fg=primary-blue}
  \setbeamercolor{itemize subitem}{fg=primary-gold}
  \setbeamercolor{alerted text}{fg=primary-gold}
  \setbeamercolor{block title}{fg=white, bg=primary-blue}
  \setbeamercolor{block body}{bg=light-bg}
  % Semantic block variants: block=definition/concept (blue), exampleblock=
  % worked example (green/positive), alertblock=key takeaway or caution (gold).
  % Keep to <=2 colored boxes per slide across ALL three variants (INV-7).
  \setbeamercolor{block title example}{fg=white, bg=positive}
  \setbeamercolor{block body example}{bg=light-bg}
  \setbeamercolor{block title alerted}{fg=white, bg=primary-gold}
  \setbeamercolor{block body alerted}{bg=light-bg}

  % Minimal footer: course code (left) + page X / N (right), no navigation clutter
  \setbeamertemplate{navigation symbols}{}
  \setbeamertemplate{footline}{%
    \color{neutral}\footnotesize\hspace{0.7em}\@coursecodetext\hfill
    \insertframenumber/\inserttotalframenumber\hspace{0.7em}\vspace{0.3em}}
}{}
\makeatother

% -----------------------------------------------------------------------------
% TikZ setup — libraries the snippet gallery uses
% -----------------------------------------------------------------------------
\usepackage{tikz}
\usetikzlibrary{arrows.meta, positioning, calc, decorations.pathreplacing,
                fit, shapes.geometric, backgrounds}

% Shared TikZ styles — snippets and hand-written diagrams can pick these up
% via `\begin{tikzpicture}[every node/.style={...}]` or by naming specific
% styles (e.g., `\node[dag-node] (X) at (0,0) {X};`).
\tikzset{
  % Node styles for causal DAGs and similar
  dag-node/.style = {
    draw=primary-blue, fill=light-bg, rounded corners=2pt,
    minimum width=1.2cm, minimum height=0.9cm, align=center, font=\small
  },
  decision-node/.style = {
    draw=primary-gold, fill=white, diamond, aspect=1.8,
    minimum width=1.8cm, minimum height=1.2cm, align=center, font=\small,
    inner sep=1pt
  },
  % Edge styles — observed vs counterfactual
  observed-edge/.style = {-{Latex[length=2.5mm]}, thick, draw=jet},
  counterfactual-edge/.style = {-{Latex[length=2.5mm]}, thick, draw=neutral, dashed},
  confound-edge/.style = {-{Latex[length=2.5mm]}, thick, draw=negative, dashed},
  % Data-point styles for plots
  observed-dot/.style = {fill=primary-blue, draw=primary-blue, circle, inner sep=1.2pt},
  counterfactual-dot/.style = {fill=white, draw=neutral, circle, inner sep=1.2pt}
}

% -----------------------------------------------------------------------------
% Convenience macros
% -----------------------------------------------------------------------------
\newcommand{\muted}[1]{\textcolor{neutral}{#1}}
\newcommand{\key}[1]{\textcolor{primary-gold}{\textbf{#1}}}
\newcommand{\good}[1]{\textcolor{positive}{#1}}
\newcommand{\bad}[1]{\textcolor{negative}{#1}}

% Standout frame macro: full-bleed dark background for section transitions.
% Beamer-only (uses \begin{frame}/\setbeamercolor/\usebeamerfont) -- do not
% call from an article-class file (e.g. Notes/<CODE>/*-notes.tex). \muted,
% \key, \good, \bad, and \coursecode above are plain xcolor/gdef and are
% safe to use from either class; verified when Notes/article-class support
% was added.
% Expand: \transitionslide{Your transition title}
\newcommand{\transitionslide}[1]{%
  \begingroup
  \setbeamercolor{background canvas}{bg=primary-blue}
  \begin{frame}[plain]
    \centering
    \vfill
    {\usebeamerfont{title}\color{white}\Large #1\par}
    \vfill
  \end{frame}
  \endgroup
}

% Section divider frame (Beamer-only), an alternative to \transitionslide with a
% darker two-line style: near-black (jet) background, a small white label line
% above a larger gold title, and a thin gold rule along the bottom edge.
% Expand: \sectiondivider{Section 2}{Pointers}
\newcommand{\sectiondivider}[2]{%
  \begingroup
  \setbeamercolor{background canvas}{bg=jet}
  \begin{frame}[plain]
    \vspace*{3.0cm}
    {\color{white}\ttfamily\large #1\par}
    \vspace{0.5cm}
    {\color{primary-gold}\LARGE #2\par}
    \begin{tikzpicture}[remember picture, overlay]
      \draw[primary-gold, line width=2.5pt]
        ($(current page.south west)+(0,0.1)$) -- ($(current page.south east)+(0,0.1)$);
    \end{tikzpicture}
  \end{frame}
  \endgroup
}

% -----------------------------------------------------------------------------
% Channel watermark — "Concepts That Click" (YouTube).
%
% Article-class files (CompetitiveExam Q/A sets, Notes, Assignments) get a
% light diagonal draft watermark on every page plus a centered footer naming
% the channel. Beamer decks get a subtle TikZ background watermark instead
% (the footline already carries the course code). Centralized here so every
% MCQ PDF is branded without per-file edits.
% -----------------------------------------------------------------------------
\makeatletter
\@ifclassloaded{article}{%
  \usepackage{draftwatermark}
  \SetWatermarkText{Concepts That Click}
  \SetWatermarkScale{1}
  \SetWatermarkFontSize{2cm}
  \SetWatermarkLightness{0.86}
  \SetWatermarkAngle{45}
  \usepackage{fancyhdr}
  \pagestyle{fancy}
  \fancyhf{}
  \fancyfoot[C]{\footnotesize\textcolor{neutral}{Concepts That Click~|~YouTube: \href{https://www.youtube.com/@concepts.thatclick}{@concepts.thatclick}}}
  \renewcommand{\headrulewidth}{0pt}
  \renewcommand{\footrulewidth}{0pt}
  \fancypagestyle{plain}{%
    \fancyhf{}%
    \fancyfoot[C]{\footnotesize\textcolor{neutral}{Concepts That Click~|~YouTube: \href{https://www.youtube.com/@concepts.thatclick}{@concepts.thatclick}}}%
    \renewcommand{\headrulewidth}{0pt}%
    \renewcommand{\footrulewidth}{0pt}%
  }
}{}
\@ifclassloaded{beamer}{%
  \setbeamertemplate{background}{%
    \begin{tikzpicture}[remember picture, overlay]
      \node[rotate=45, opacity=0.055, align=center]
        at (current page.center)
        {\Huge\textcolor{neutral}{Concepts That Click}};
    \end{tikzpicture}%
  }
}{}
\makeatother 
\coursecode{JKSSB}
\renewcommand{\thesection}{47.\arabic{section}}
\newtheorem{example}{Example}[section]
\newenvironment{definitionbox}[1]{%
  \par\noindent\textbf{Definition (#1).}\ \itshape
}{\par\vspace{0.3em}}
\title{Humanware and IT Governance --- Detailed Lecture Notes}
\author{JKSSB Computer Awareness}
\date{\today}

\begin{document}
\maketitle

\section{What humanware means}
A computer system is not only metal and programs. It also needs people. The
human element of a computer system is called \textbf{humanware}, and it is
also known as \textbf{liveware} or \textbf{peopleware}. Humanware means the
people who design, build, operate, manage and use computers. It is
\emph{neither hardware nor software}; it is the people component of the
system.

\begin{definitionbox}{Humanware}
The people involved with a computer system --- its users, operators,
administrators, programmers and support staff. Also called liveware or
peopleware.
\end{definitionbox}

\begin{center}
\begin{tabular}{p{0.24\textwidth}p{0.66\textwidth}}
\toprule
\textbf{Term} & \textbf{Meaning in this lesson} \\
\midrule
Hardware & The physical parts of the computer. \\
Software & The programs and instructions the computer runs. \\
Humanware & The people who use and manage the hardware and software. \\
User & The person who uses the finished system to do a task. \\
Administrator & The person who manages the system, network or database. \\
\bottomrule
\end{tabular}
\end{center}

The rest of this lesson uses these terms in one consistent sense.
\textbf{Hardware} is the machinery. \textbf{Software} is the set of
instructions. \textbf{Humanware} is the people.

\section{The people in a computer system}
Humanware is not a single role. Many different people take part, and each has
a distinct job. The main categories are the following.

\begin{itemize}
  \item \textbf{User (or end user)} --- the person who uses the finished
    system to perform a task. A clerk entering data, a student reading a
    portal and a customer paying a bill are all users.
  \item \textbf{Operator} --- the person who runs the computer system, feeds
    jobs into it and monitors routine operation.
  \item \textbf{Administrator} --- the person who manages the system. A
    \textbf{system administrator} manages servers, user accounts and access
    rights; a \textbf{network administrator} manages the network; a
    \textbf{database administrator} manages the database.
  \item \textbf{Programmer (developer)} --- the person who writes, tests and
    maintains the software.
  \item \textbf{IT support (help desk) and technician} --- the people who
    answer user problems, troubleshoot faults, and install, repair and
    maintain the hardware.
\end{itemize}

The distinction to hold is simple. Administrators \emph{manage} the system,
operators \emph{run} it, and users \emph{use} its output.

\section{Hardware, software and humanware}
A working computer system needs all three components. Hardware is the
physical machine. Software is the set of programs that tell the hardware what
to do. Humanware is the people who design, run and use both. Hardware without
software does nothing, and both are useless without the people who operate
them. This is why humanware is taught as a full component, not as an
afterthought.

\begin{example}
A government office installs new computers (hardware) and an accounting
program (software). The system still does not work by itself. A programmer
must maintain the program, a system administrator must create user accounts,
an operator must run the daily jobs, and the office staff (users) must enter
the data. Those people are the humanware.
\end{example}

\section{The role of IT in governance}
\textbf{Governance} is the process of decision-making and its implementation
in an organisation or a state. \textbf{IT} (Information Technology) and
\textbf{ICT} (Information and Communication Technology) let government carry
out that work over computers and networks. The role of IT in governance is to
make government faster, cheaper, more transparent and more accessible to the
people it serves.

The main benefits of using IT in governance are the following.

\begin{itemize}
  \item \textbf{Transparency} --- citizens can see the rules, the status of
    their applications and the decisions taken.
  \item \textbf{Accountability} --- actions are recorded and can be traced.
  \item \textbf{Efficiency} --- less paperwork and faster processing.
  \item \textbf{Convenience} --- services are available from home and around
    the clock.
  \item \textbf{Lower cost} and \textbf{reduced corruption}, because manual
    handling and physical counters are replaced by digital processes.
\end{itemize}

The overall aim is \textbf{citizen-centric} service: the citizen is placed at
the centre of every process.

\section{IT governance versus e-governance}
These two terms are a common near-neighbour pair, and the difference matters.

\begin{itemize}
  \item \textbf{IT governance} is the process by which an organisation
    directs and controls its own IT resources so that they support its goals.
    It is \emph{governance of IT}.
  \item \textbf{E-governance} is the use of IT by government to govern and
    deliver public services. It is \emph{IT in governance}.
\end{itemize}

The trap is to swap them. IT governance is about managing IT well inside an
organisation; e-governance is about using IT to govern the state and serve
citizens.

\section{E-governance basics}
\textbf{E-governance} is the application of ICT by government to deliver
services, share information and transact with citizens, businesses and other
government agencies. It uses the internet, databases and digital records in
place of paper and physical counters. It makes government services faster,
wider-reaching and easier to audit.

\begin{definitionbox}{E-governance}
The use of information and communication technology to improve the way
government serves and interacts with the public.
\end{definitionbox}

E-governance is usually described by the parties who interact. The four
standard types are the following.

\begin{center}
\begin{tabular}{p{0.14\textwidth}p{0.32\textwidth}p{0.44\textwidth}}
\toprule
\textbf{Type} & \textbf{Full form} & \textbf{Interaction} \\
\midrule
G2C & Government to Citizen & Services to the general public. \\
G2G & Government to Government & Between government departments. \\
G2B & Government to Business & With companies and traders. \\
G2E & Government to Employee & With government staff. \\
\bottomrule
\end{tabular}
\end{center}

This lesson keeps the treatment at awareness level. Named platforms and
detailed G2x matching belong to Lesson 57.

\section{Digital public services}
A \textbf{digital public service} is a government service that a citizen can
obtain online. Common examples are certificates, licences, land records, bill
payment, applications and grievance redressal. Citizens reach these services
through service portals and through \textbf{Common Service Centres (CSCs)},
which provide a physical access point where help is needed. Typical features
are online application, status tracking, digital payment and downloadable
documents. A service counts as digital when the citizen can complete the task
without visiting an office in person.

\section{Electronic records}
An \textbf{electronic record} is information created, stored and retrieved in
electronic (digital) form. Examples include e-documents, e-files, scanned
papers, database entries and email. Electronic records are easy to store,
search, retrieve, share and back up, and they replace large volumes of paper.
An electronic record can have legal validity when it is authenticated, for
example with a digital signature.

\begin{definitionbox}{Electronic record}
Information that is created, stored and retrieved in electronic form, such as
an e-document, an e-file or a database entry.
\end{definitionbox}

\section{Authentication}
\textbf{Authentication} is the process of verifying that a user, device or
system is who or what it claims to be. It answers the question: ``Who are
you?'' Common methods are a password, a PIN, an \textbf{OTP} (One-Time
Password), a biometric such as a fingerprint or iris, and a digital
certificate. Authentication must succeed before the system grants access.

\begin{definitionbox}{Authentication}
The process of confirming a claimed identity before access to a system is
allowed.
\end{definitionbox}

Authentication is often confused with \textbf{authorization}.
\textbf{Authorization} is the process of granting or denying permission ---
it decides what a verified identity is allowed to do. Authentication proves
identity; authorization grants or refuses access rights. Authentication
always comes first.

Authentication relies on \emph{factors}. There are three standard types:
something you \textbf{know} (a password or PIN), something you \textbf{have}
(an OTP on your phone, a token or a smart card), and something you \textbf{are}
(a fingerprint, iris or face). When two different factors are combined, it is
called \textbf{two-factor authentication (2FA)}; when two or more are used, it
is \textbf{multi-factor authentication (MFA)}. Two passwords are not two
factors, because they are the same factor used twice.

\begin{example}
A candidate logs in to a government portal with a password and is then asked
for an OTP sent to the registered mobile number. The password is something the
candidate knows; the OTP is something the candidate has. Together they form
two-factor authentication. The portal then decides which forms the candidate
may open --- that decision is authorization, not authentication.
\end{example}

\section{Exam retrieval and common confusions}
Five distinctions prevent most errors on this topic.
\begin{enumerate}
  \item Humanware means \textbf{people}. It is neither hardware nor software;
    do not confuse it with firmware (which is software stored in hardware).
  \item \textbf{IT governance} is governance \emph{of} IT;
    \textbf{e-governance} is IT \emph{in} governance. Do not swap them.
  \item \textbf{G2C} is Government to \textbf{Citizen}, not Government to
    Company. The business interaction is G2B.
  \item \textbf{Authentication} proves identity; \textbf{authorization}
    grants permission. Authentication comes first.
  \item An \textbf{electronic record} is digital information. A printed paper
    file is a physical record, not an electronic one.
\end{enumerate}

\begin{example}
An item asks which component of a computer system consists of the people who
operate and manage it. Hardware is the machinery and software is the set of
instructions, so both are wrong. The correct answer is humanware. If the stem
instead asks what improves transparency and accountability in government, the
answer is the role of IT in governance, and if it asks what verifies a user's
identity, the answer is authentication.
\end{example}

\subsection*{Exercises}
\begin{enumerate}[label=\arabic*.]
  \item Define humanware and name four categories of people who form it.
  \item Distinguish hardware, software and humanware, and explain why all
    three are needed.
  \item Give three benefits of using IT in governance.
  \item Distinguish IT governance from e-governance.
  \item What are the four types of e-governance interaction, and what does
    each connect?
  \item Define an electronic record and give two examples.
  \item Distinguish authentication from authorization, and name the three
    standard authentication factors.
\end{enumerate}

\subsection*{Solutions}
\begingroup\small
\begin{enumerate}[label=\arabic*.]
  \item Humanware (also called liveware or peopleware) is the human element of
    a computer system. Four categories are users (end users), operators,
    administrators, and programmers (with IT support and technicians).
  \item Hardware is the physical machine, software is the set of programs,
    and humanware is the people. Hardware without software does nothing, and
    both need people to design, run and use them, so all three are required.
  \item Any three of: transparency, accountability, efficiency, convenience,
    lower cost, reduced corruption, and citizen-centric service.
  \item IT governance is the direction and control of an organisation's own
    IT resources --- governance of IT. E-governance is the use of IT by
    government to deliver public services --- IT in governance.
  \item G2C (Government to Citizen) serves the public; G2G (Government to
    Government) links departments; G2B (Government to Business) links
    companies; G2E (Government to Employee) links government staff.
  \item An electronic record is information created, stored and retrieved in
    electronic form. Examples include an e-document, an e-file, a scanned
    paper or a database entry.
  \item Authentication verifies identity (``who are you?''); authorization
    grants or denies permission (``what may you do?''). The three factors are
    something you know, something you have, and something you are.
\end{enumerate}
\endgroup
\end{document}
